% core_resume_template_v3.tex - ATS resume, plan v3 (2026-10-17).
%
% Single column, standard section headers on their own line, no columns or
% tables (Textkernel quality codes 151/433), <= 2 pages (German norm).
% Font recipe validated on this repo's tectonic (2026-10-17 test): cfr-lm under
% XeTeX dropped U+2014/U+2019 ("Software Engineer Soham Infotech", "Runners Up");
% Pagella alone merged words in extraction ("efficientoffice"). Only Pagella +
% Ligatures=NoCommon + \raggedright + \spaceskip 0.34em extracts cleanly.
% Every token is pre-rendered by templates/render.py (fill.py never escapes).
\documentclass[a4paper,10pt]{article}

\usepackage[left=1.5cm,right=1.5cm,top=1.2cm,bottom=1.2cm]{geometry}
\usepackage{fontspec}
\setmainfont{texgyrepagella}[Ligatures=NoCommon, Extension=.otf, UprightFont=*-regular, BoldFont=*-bold, ItalicFont=*-italic, BoldItalicFont=*-bolditalic]
\usepackage{enumitem}
\usepackage{titlesec}
\usepackage[hidelinks]{hyperref}
\hypersetup{pdftitle={Meetkumar Riteshkumar Raval - CV}, pdfauthor={Meetkumar Riteshkumar Raval}}

\pagestyle{empty}
\raggedright
\setlength{\parindent}{0pt}
\setlength{\parskip}{0pt}
\setlength{\spaceskip}{0.34em plus 0.08em minus 0.02em}
\setlength{\xspaceskip}{0.34em plus 0.08em minus 0.02em}

\setlist[itemize]{leftmargin=1.2em,nosep,topsep=1pt,itemsep=1pt,parsep=0pt}
\titleformat{\section}{\large\bfseries}{}{0pt}{}[\vspace{-5pt}\rule{\textwidth}{0.5pt}]
\titlespacing*{\section}{0pt}{8pt}{4pt}

\begin{document}

\begin{center}
  {\LARGE\bfseries Meetkumar Riteshkumar Raval}\par\vspace{3pt}
  {\large Werkstudent AI \& Python Engineering | MSc Artificial Intelligence}\par\vspace{3pt}
  {\small Essen, Germany | +49 15565 468731 | \href{mailto:ravalmeett@gmail.com}{ravalmeett@gmail.com} | \href{https://www.linkedin.com/in/meetraval/}{linkedin.com/in/meetraval} | \href{https://github.com/ravalmeet}{github.com/ravalmeet} | \href{https://raval.vercel.app/}{raval.vercel.app}}\par\vspace{2pt}
  {\small Available immediately | 20 h/week during lecture periods, up to 40 h/week in semester breaks or internships | Student residence permit valid until December 2028 - eligible for working-student jobs and internships in Germany}\par
\end{center}

\section*{Profile}
MSc Artificial Intelligence student at BTU Cottbus-Senftenberg and former Software Engineer with hands-on work in LLMs and AI. Proficient in designing automated pipelines and AI agents for data transformation and code generation. Role-relevant skills: Python backend services, PyTorch, LangChain and Docker.\par

\section*{Technical Skills}
\begin{itemize}
\item \textbf{Programming Languages:} Python backend services, Python, C++, JavaScript, TypeScript, HTML
\item \textbf{AI \& Machine Learning:} Large Language Models (LLMs), PyTorch, LangChain, Artificial Intelligence (AI), Computer Vision, Deep Learning, Machine Learning (ML)
\item \textbf{Data \& Analytics:} MongoDB
\item \textbf{Software, Cloud \& DevOps:} Docker, FastAPI, Django, Automated Workflows
\item \textbf{Tools \& Methods:} Git, Agile Development
\end{itemize}

\section*{Work Experience}
\textbf{Software Engineer} | Soham Infotech | Gujarat, India \hfill Dec 2023 - Jan 2025\par
\begin{itemize}
\item Developed modular API endpoints and quantitative functions that increased profitability of financial decision-making processes.
\item Designed and built automated data transformation pipelines and backend components, integrating AI to support real-time trading operations.
\item Collaborated within an agile engineering team, utilizing Git and GitHub while applying software engineering best practices for high-reliability systems.
\item Focused on black-box and gray-box models, achieving benchmark-surpassing performance across evaluated metrics.
\item \textbf{Key skills:} Artificial Intelligence (AI), Python, Deep Learning, Quantitative Modeling, Automated Workflows, Software Engineering Best Practices, Agile Development, Git
\end{itemize}
\vspace{3pt}
\textbf{Machine Learning Intern} | Pantechelearning | Chennai, India \hfill Dec 2021 - Feb 2022\par
\begin{itemize}
\item Implemented AI-driven machine learning models for diverse projects, such as text-to-speech applications and high-functional visual-activity analysis.
\item Developed automated data extraction and preprocessing pipelines, guaranteeing accurate AI model training and evaluation inputs.
\item \textbf{Key skills:} Artificial Intelligence (AI), Python, Machine Learning (ML), Deep Learning, Natural Language Processing (NLP), Computer Vision, Data Preprocessing, Text-to-Speech, Django
\end{itemize}
\vspace{3pt}

\section*{Projects}
\textbf{High Frequency Trading Platform}\par
\begin{itemize}
\item Developed AI-driven real-time data scraping, analysis, and filtering pipeline for high-frequency trading decisions.
\item Implemented AI models combined with quantitative logic to execute market actions and manage unexpected events.
\item \textbf{Tech:} Artificial Intelligence (AI), Python, Mathematics, Quantitative Modeling, Django, Data Analytics, Machine Learning (ML), Deep Learning, Real-Time Systems, Data Scraping
\end{itemize}
\vspace{3pt}
\textbf{Automated Resume Generation} | \href{https://job-automation-raval.streamlit.app/}{job-automation-raval.streamlit.app}\par
\begin{itemize}
\item Created LLM-based automated transformation system that converts code and text into structured documents using prompt engineering.
\item Engineered agent workflows with Python and Transformers to categorize roles and generate resumes via automated pipelines.
\item \textbf{Tech:} Large Language Models (LLMs), Python, Streamlit, Natural Language Processing (NLP), Transformers, AI Agents, Automated Workflows, Git
\end{itemize}
\vspace{3pt}

\section*{Education}
\textbf{MSc Artificial Intelligence} \hfill Apr 2025 - Oct 2027 (expected)\par
Brandenburg University of Technology Cottbus-Senftenberg (BTU) | Germany\par
Faculty 1 - Mathematics, Computer Science, Physics, Electrical Engineering and Information Technology\par
Focus areas: Core AI Fundamentals, Mathematics, Problem Solving, Critical Thinking\par
\vspace{3pt}
\textbf{BE Computer Engineering} \hfill Jun 2020 - May 2024\par
Gujarat Technological University | India\par
\vspace{3pt}

\section*{Research}
\begin{itemize}
\item \textbf{An Empirical Study on Large Language Models for Sarcasm Detection in Low-resourced Indic Languages}\\ Conducting research on NLP models aimed at detecting and analyzing sarcasm across various Indian regional dialects.
\end{itemize}

\section*{Achievements}
\begin{itemize}
\item Smart India Hackathon - 2021: Runner's Up - State Level
\item Smart India Hackathon - 2022: Runner's Up - National Level
\item Azadi ka Amrit Mahotsav - 2023: Runner's Up - State Level
\item AI Guild Hackathon - 2025: Winner - Berlin and Brandenburg State, Germany
\end{itemize}

\section*{Certifications}
\begin{itemize}
\item Google - Big Data and Machine Learning
\item Google - Machine Learning with TensorFlow
\item How Google does Machine Learning
\item Google - Machine Learning
\item Machine Learning with Python
\item Data Science Introduction
\item Tools for Data Science
\item Problem Solving (Intermediate) - HackerRank
\item Java Programming: Complete Beginner to Advanced - Udemy
\item Web Development Masterclass Course - Udemy
\end{itemize}

\section*{Additional Projects}
\begin{itemize}
\item TradingView-Data Fetch: Individual project; open source on GitHub with 40+ stars and 24+ forks | \href{https://github.com/ravalmeet/TradingView}{github.com/ravalmeet/TradingView}
\item Freelance web projects delivered for clients: Nirvana Food Management, Select-Your-Tutor, Shivshivay Himalaya (Shivshivay.com)
\end{itemize}

\section*{Languages}
English - Fluent (C2) | German - Basic (A2), currently learning | Hindi - Native | Gujarati - Native\par

\end{document}
